\section{Troubleshooting}\label{sec:troubleshooting}

\begin{description}[style=nextline, leftmargin=0pt, labelwidth=0pt, labelsep=0pt]
  \item[No image or a stale image] Check the source signal, cable, input
  gain, and Run light. Ensure LO Freq is below HI Freq. If you loaded a
  frozen patch, resume capture to acquire new history.
  \item[Bands are too broad or changes blur together] Reduce Smooth and
  Average, and compare window functions. Spectre's FFT length is fixed;
  use Fourier if you need a longer frame to separate nearby frequencies.
  \item[Levels differ from another analyzer] Match input gain, window,
  FFT length, smoothing, and amplitude reference. Set Slope to zero and
  account for summed polyphonic voices. Spectre starts at $+6\,\mathrm{dB}$
  input gain; Fourier starts at $0\,\mathrm{dB}$. A spectral peak is not a
  broadband RMS measurement.
  \item[Large regions reach the same end color] In Decibels mode, raise
  the upper handle to distinguish strong values in retained history. Adjust
  the lower handle for quiet detail. Choose Linear for less visible weak background
  detail, or raise Floor in Decibels mode.
  Color saturation does not mean the stored magnitude was clipped.
  \item[Old and new bands disagree after a sample-rate change] Allow a full
  sweep to replace history acquired at the old sample rate before comparing
  frequencies. Check the frequency bounds as well.
  \item[Analysis uses too much CPU] Spectre performs analysis on Rack's
  audio engine thread despite having no audio outputs. Turn Run off when
  inspecting a frozen image to pause buffering and FFT processing. Its FFT
  length and hop size are fixed; Fourier provides adjustable settings when
  you need to trade frequency detail or update rate for analysis workload.
\end{description}
